\clearpage


\chapter{Технологическая часть}\label{ch:-2}

В этом разделе обоснованы средства реализации, а так же представлены реализации алгоритмов и функциональные тесты.


\section{Средства реализации}\label{sec:-2}

Для реализации алгоритмов в этой лабораторной работы был выбран язык $Python^{[1]}$ по следующим причинам:
\begin{itemize}
    \item python предоставляет необходимые структуры данных, такие как списки (массивы);
    \item поддерживаемые библиотеки, такие как matplotlib позволяют визуализировать результаты исследования.
\end{itemize}

\clearpage


\section{Реализация алгоритмов}

В листингах \ref{lst:lsearch} --- \ref{lst:bsearch} представлены реализации алгоритмов.

\begin{lstlisting}[style=python, label=lst:lsearch,caption=Алгоритм использующий полный перебор]
def search_linear(arr, elem):
    length = len(arr)

    if length == 0:
        return -1, 0

    cmp = 0

    for i in range(length):
        cmp += 1

        if arr[i] == elem:
            return i, cmp

    return -1, cmp
\end{lstlisting}
%\clearpage


\begin{lstlisting}[style=python, label=lst:bsearch,caption=Алгоритм использующий бинарный поиск]
def search_binary(arr, elem):
    length = len(arr)

    if length == 0:
        return -1, 0

    cmp = 0
    left = 0
    right = len(arr) - 1
    while left <= right:
        cmp += 1
        mid = (left + right) // 2

        if arr[mid] == elem:
            return mid, cmp
        elif arr[mid] < elem:
            left = mid + 1
        else:
            right = mid - 1

    return -1, cmp

\end{lstlisting}
\clearpage


\section{Функциональные тесты}
В таблице \ref{tbl:func_test} приведены функциональные тесты, на которых тестировалась программа.
\begin{table}[h]
    \begin{center}
        \caption{\label{tbl:func_test} Функциональные тесты, на которых были протестированы алгоритмы (в результате -- индекс элемента).}
        \begin{tabular}{|c|c|c|c|}
            \hline
            Массив & Искомый элемент & Линейный поиск & Бинарный поиск
            \\ \hline
            [ 1, 2, 3, 4, 5 ] & 1 & 0 & 0
            \\ \hline
            [ 1, 2, 3, 4, 5 ] & 5 & 4 & 4
            \\ \hline
            [ 1, 2, 3, 4, 5 ] & 3 & 2 & 2
            \\ \hline
            [] & 3 & -1 & -1
            \\ \hline
            [ 1, 2, 3, 4, 5 ] & 6 & -1 & -1
            \\ \hline
        \end{tabular}
    \end{center}
\end{table}

\section*{Вывод}

В технологической части были представлены листинги реализованных алгоритмов, включая линейный и бинарный поиск по массиву.
Также были проведены функциональные тесты, которые подтвердили корректность работы алгоритмов на различных входных данных, включая граничные случаи.
Результаты тестов показали, что разработанные реализации соответствуют поставленным требованиям.

